\documentclass{article}
\usepackage{amsmath,amssymb,amsthm}
\newtheorem{theorem}{Theorem}
\newtheorem{lemma}{Lemma}
\newtheorem{corollary}{Corollary}
\title{GST new native infinity controller: quantized approximate Betti matching plus bounded Chow complexity forces an ACTUAL original-X algebraic source}
\date{12 November 2026}
\begin{document}\maketitle

\section{Precise universal prescribed-target problem}
For arbitrary smooth projective complex original X, codimension p
and prescribed rational Hodge class alpha, construct an ACTUAL rational
Chow cycle on the SAME X with CANONICAL rational Betti class alpha.
The owner's original GST uses its genuine twelve-state U,J,M,[J,M]
wave fold, physical Navigation, N-cohomology and cardinally indexed
topologies. The previous bounded-complexity theorem reduced the older
aleph-one condition to a COUNTABLE sequence of exactly target-matching
genuine algebraic cycles of controlled degree and denominators.

The present theorem REMOVES exact target matching from the input:
only sufficiently CLOSE rational Betti classes are required.
The proof is based on INTEGRAL QUANTIZATION, so no numerical
continuity can be substituted for the actual rational class lattice.

\section{An exact sharp rational-lattice threshold}
Let L be ANY finite-rank free abelian group, let V=L tensor_Z Q,
and choose an INTEGRAL basis e_1,...,e_r of L.
Use the maximum coordinate norm on V_R.
Fix alpha in V with a positive integer d such that d alpha in L.
Let beta=(1/m)lambda with lambda in L and 1<=m<=C.
No geometry is assumed in this exact arithmetic lemma.

\begin{lemma}[Sharp denominator-controlled exactness]
If
\[
\boxed{\|\beta-\alpha\|_\infty<1/(d C),}
\]
then beta=alpha EXACTLY in V.
Equivalently, any distinct rational beta of denominator
at most C is separated from alpha by at least 1/(d C).
\end{lemma}
\begin{proof}
The vector d m (beta-alpha)=d lambda-m(d alpha)
is an INTEGRAL lattice vector. If beta is not alpha,
some integral coordinate of this vector has absolute
value at least 1. That coordinate of beta-alpha has
absolute value at least 1/(d m)>=1/(d C).
This contradicts the strict upper bound.
\end{proof}
The threshold is SHARP as a universal statement:
with L=Z, alpha=2/5, d=5, C=7 and beta=3/7,
\[
|\beta-\alpha|=|3/7-2/5|=1/35=1/(dC),
\]
so the strict inequality cannot be replaced by a
nonstrict inequality without additional hypotheses.
This example is a pure exact arithmetic diagnostic,
not a special-family substitute for original-X geometry.

\section{The new arbitrary-original-X native approximate-source theorem}
Let pi:calX->B be smooth PROJECTIVE over an irreducible smooth
complex algebraic CURVE. Let b0 be the original fiber X, and
let alpha_b be a GLOBAL FLAT rational cohomology section of
degree 2p, with alpha_b0 the prescribed original rational Hodge
class. Work in a simply connected analytic neighborhood
of b0, identifying the FREE part of integral cohomology
L_b=H^{2p}(X_b,Z)/torsion across all fibers by Gauss--Manin.
Choose an integral basis of that local system and a
positive integer d clearing the denominators of alpha.
Fix a relative projective embedding.

For each b_n tending to b0, suppose a GST native geometry
recipe produces an ACTUAL rational Chow difference
\[
Z_n=(P_n-Q_n)/m_n\in CH^p(X_{b_n})_Q
\]
from genuine effective algebraic cycles on the original
moving fiber, with
\[
1\le m_n\le C,\qquad
\deg(P_n),\deg(Q_n)\le D
\]
for fixed positive integers C,D independent of n.
The new hypothesis is NOT that cl(Z_n)=alpha_b_n.
It is only the quantitative inequality
\[
\boxed{
\|\mathrm{cl}^{can}_{X_{b_n}}(Z_n)-\alpha_{b_n}\|_\infty
<1/(dC)
}
\tag{1}
\]
at every n, with the canonical rational Betti class and
the fixed integral cohomology coordinates.

\begin{theorem}[Quantized approximate GST infinity-return]
Under the above hypotheses the ORIGINAL X has a
genuine rational algebraic cycle
\[
\boxed{Z_\alpha\in CH^p(X)_Q,\quad
\mathrm{cl}^{can}_X(Z_\alpha)=\alpha.}
\]
No original-X Chow preimage, no uncountably many exact
source-bearing fibers, no generic-fiber algebraic
representative and no equality cl(Z_n)=alpha_b_n
is assumed as input.
\end{theorem}
\begin{proof}
Each m_n cl(Z_n) is the canonical class of an
INTEGRAL algebraic cycle and is therefore in the
integral Betti cohomology lattice L_b_n
(modulo torsion). The globally flat rational alpha
has the SAME denominator d after identifying the
local integral lattices over a simply connected
neighborhood of the original base point.
The sharp rational-lattice lemma therefore applies
to EACH fiber and inequality (1) and gives
\[
\mathrm{cl}_{X_{b_n}}(Z_n)=\alpha_{b_n}
\]
EXACTLY, not asymptotically.
Now the independently proved bounded-complexity
relative Chow specialization theorem applies:
the finite profiles $(m_n,\deg P_n,\deg Q_n)$
have a fixed triple on an infinite subsequence
$b_{n_k}\to b0$. The good relative Chow pair
parameter is proper over B; its cohomological
label is locally constant in Gauss--Manin;
hence the good source locus is closed, contains
the infinite subsequence and contains b0.
An ACTUAL pair $(P0,Q0)$ over original X yields
$Z_\alpha=(P0-Q0)/m$ with canonical class alpha.
\end{proof}

\begin{corollary}[Exact approximate-source obstruction to a missing Chow class]
If the prescribed original rational Hodge alpha
has no algebraic representative, then for ANY
sequence b_n->b0 of genuine algebraic fiber
cycles with uniformly bounded denominator and
geometric degree, their canonical Betti classes
CANNOT eventually enter the target ball
of radius 1/(dC).
If they converge to alpha in the canonical
rational cohomology coordinates, at least one
of the denominator or the effective degrees
MUST escape every finite bound.
\end{corollary}

\section{A usable GST native acceptance test}
This is a STRICT improvement in the source-input condition
over the prior exact-match aleph-one and countable-fiber theorems:
the original GST F-wave may construct sources whose
canonical cohomology is only APPROXIMATE to the marked
alpha. It need not solve the equality at every step.
If its actual algebraic output on moving projective
fibers simultaneously has a bounded rational denominator,
bounded genuine algebraic degrees, and enters the explicit
quantized target ball, the output on original X is
PROVED by the preceding theorem.

Crucially, the present mathematical results do not
establish that the owner's currently defined native U/J/M
curvature, cardinal-world topology or infinite iteration
AUTOMATICALLY produces ANY such genuine cycles on
arbitrary original X.
The physical ON fold's normalized stable projector
loses a source direction and the OFF fold retains
Jordan growth, as independently proved in the prior
native complexity derivation. Native convergence
of abstract wave coordinates is NOT the same as
canonical convergence of ACTUAL algebraic cycle classes.

This theorem is a noncircular CONSTRUCTIVE conditional
criterion for source birth on original X from quantitative
data on other fibers; it does not prove the full universal
Hodge conjecture.

\section{Sharp infinity-control cost: inverse-error denominator growth}
There is a stronger EXACT quantified consequence
which remains valid even when the actual Chow
degrees are NOT bounded.

Fix the rational original target alpha with
a clearing denominator d, and a true
moving-fiber algebraic rational cycle
beta_n with a chosen denominator m_n
in a common integral Gauss-Manin lattice.
Suppose beta_n is NOT exactly alpha_b_n.
The same integral-lattice argument gives
\[
\boxed{
\|\beta_n-\alpha_{b_n}\|_\infty
\ \ge\ \frac1{d\,m_n}.
}
\]
Equivalently, if a GST process produces
a nonmatching algebraic cycle with
Betti error strictly below epsilon_n>0,
then it MUST have
\[
\boxed{
m_n>\frac1{d\,\epsilon_n}.
}
\]
For exponential error targets
$\epsilon_n=2^{-n}$,
every continued NONMATCHING native
approximation requires
$m_n>2^n/d$.
This is an exact ALL-n theorem, not a
computational trend.
It gives the desired quantitative
infinity-control dichotomy:

Either a source matches the marked
rational Hodge class EXACTLY in finite time,
or the denominators of endlessly improving
nonmatching algebraic approximants
escape at least reciprocally with the
error bound. This arithmetic fact
does not manufacture those approximants.

The bound is sharp without more input,
as witnessed by the exact rational
one-dimensional data alpha=2/5 and
beta=3/7, with d=5,m=7 and
absolute error precisely 1/35.

\end{document}
